\documentclass[conference]{IEEEtran}
\usepackage{cite}
\usepackage{amsmath,amssymb,amsfonts}
\usepackage{algorithmic}
\usepackage{graphicx}
\usepackage{textcomp}
\usepackage{xcolor}
\usepackage{booktabs}
\usepackage{multirow}
\usepackage{listings}
\usepackage{hyperref}

\lstset{
    basicstyle=\ttfamily\footnotesize,
    breaklines=true,
    frame=single,
    language=C,
    numbers=left,
    numberstyle=\tiny,
    tabsize=2
}

\begin{document}

\title{HW5: Domain-Specific
Accelerator}

\author{\IEEEauthorblockN{Cheng Hao Lai, 112550171}}

\maketitle

\begin{abstract}
This report presents an incremental hardware-software co-design approach for accelerating a Convolutional Neural Network (CNN) inference on the RISC-V Aquila System-on-Chip (SoC). Starting from a pure software baseline implementation that requires 54,523 ms for handwritten digit recognition, we progressively apply various optimization techniques including MMIO-based Domain-Specific Accelerator (DSA), direct gather, weight preloading, and pooling acceleration. Our final optimized implementation achieves a speedup of 12.72$\times$ while maintaining 98\% classification accuracy on the MNIST-like test dataset.
\end{abstract}

\begin{IEEEkeywords}
CNN acceleration, RISC-V, FPGA, hardware-software co-design, DSA, Aquila SoC
\end{IEEEkeywords}

\section{Introduction}

Convolutional Neural Networks (CNNs) have become fundamental building blocks in modern machine learning applications. However, deploying CNN inference on resource-constrained embedded systems remains challenging due to the computational intensity of convolution operations. This work explores hardware-software co-design techniques to accelerate CNN inference on the Aquila RISC-V SoC.

The target application is a 6-layer CNN model for handwritten digit recognition, consisting of two convolutional layers, two average pooling layers, and two fully connected layers. The baseline pure-software implementation requires approximately 54.5 seconds to classify 100 test images, which is impractical for real-time applications.

We present an incremental optimization methodology, progressively identifying and addressing performance bottlenecks through a combination of hardware acceleration and software optimization. Each optimization is controlled by a compile-time flag in the Makefile, allowing systematic evaluation.

\section{System Architecture}


\subsection{CNN Model Architecture}

The CNN model used for evaluation consists of the following layers:

\begin{enumerate}
    \item \textbf{Conv1}: 1 input channel, 3 output channels, 5$\times$5 kernel
    \item \textbf{Pool1}: 2$\times$2 average pooling
    \item \textbf{Conv2}: 3 input channels, 12 output channels, 5$\times$5 kernel
    \item \textbf{Pool2}: 2$\times$2 average pooling
    \item \textbf{FC1}: 192 $\rightarrow$ 30 fully connected
    \item \textbf{FC2}: 30 $\rightarrow$ 10 fully connected with ReLU
\end{enumerate}

\subsection{DSA Hardware Architecture}

The Domain-Specific Accelerator (DSA) is designed as an MMIO-mapped peripheral with the base address at \texttt{0xC4000000}. The hardware implements a floating-point Fused Multiply-Add (FMA) unit using Xilinx IP cores. The register map is shown in Table~\ref{tab:mmio_reg}.

\begin{table}[htbp]
\caption{DSA MMIO Register Map}
\label{tab:mmio_reg}
\centering
\begin{tabular}{@{}lll@{}}
\toprule
Offset & Name & Description \\
\midrule
0x000 & CTRL & Control/Status register \\
0x004 & LEN & Vector length register \\
0x008 & W\_OFF & Weight offset in BRAM \\
0x010 & IN\_A & Input buffer A (252 words) \\
0x400 & IN\_B & Weight buffer B (1024 words) \\
0x1000 & OUT & Output result (32-bit float) \\
\bottomrule
\end{tabular}
\end{table}


\subsection{MMIO Communication Mechanism}

Memory-Mapped I/O (MMIO) is the fundamental communication mechanism between the RISC-V processor and the DSA hardware. In the Aquila SoC, the address range \texttt{0xC400\_0000} to \texttt{0xC4FF\_FFFF} is allocated to the DSA device.

\subsubsection{Address Decoding}
When the processor executes a \texttt{load} or \texttt{store} instruction targeting an address within the DSA range, the \texttt{soc\_top.v} module performs address decoding:

\begin{lstlisting}[language=Verilog,caption=Address Decoding in soc\_top.v]
assign dsa_sel = (dev_addr[31:24] == 8'hC4);
assign dev_dout = (dsa_sel) ? dsa_dout : ...;
assign dev_ready = (dsa_sel) ? dsa_ready : ...;
\end{lstlisting}

When \texttt{dsa\_sel} is asserted, the device bus signals (\texttt{dev\_strobe}, \texttt{dev\_we}, \texttt{dev\_addr}, \texttt{dev\_din}) are routed to the \texttt{dsa\_feeder} module.

\subsubsection{Data Transfer Protocol}
The MMIO write sequence for a dot product operation consists of:

\begin{enumerate}
    \item \textbf{Write LEN register}: Set the vector length $N$
    \item \textbf{Write IN\_A buffer}: Write $N$ input values (activations)
    \item \textbf{Write IN\_B buffer}: Write $N$ weight values
    \item \textbf{Write CTRL register}: Set START bit to trigger computation
    \item \textbf{Poll CTRL register}: Wait for DONE bit
    \item \textbf{Read OUT register}: Retrieve the result
\end{enumerate}

Each MMIO write requires the processor to:
\begin{itemize}
    \item Generate the target address
    \item Execute the \texttt{sw} (store word) instruction
    \item Wait for the memory bus transaction to complete
\end{itemize}

At the Aquila SoC clock frequency of approximately 41 MHz, each MMIO write takes 3-5 clock cycles, resulting in significant overhead when transferring large data vectors.




\section{Optimization Methodology}

We adopt an incremental optimization approach, applying one technique at a time to clearly measure the contribution of each optimization. Each optimization is controlled by a compile-time flag, as shown in Table~\ref{tab:flags}. Table~\ref{tab:results} summarizes the performance results at each stage.

\begin{table}[htbp]
\caption{Compile-time Optimization Flags}
\label{tab:flags}
\centering
\begin{tabular}{@{}ll@{}}
\toprule
Flag & Description \\
\midrule
\texttt{USE\_HW\_DOT} & FC layer dot-product acceleration \\
\texttt{USE\_HW\_CONV3D} & Conv layer hardware acceleration \\
\texttt{USE\_HW\_DIRECT\_GATHER} & Eliminate flatten overhead \\
\texttt{USE\_HW\_CONV\_SIMPLIFIED} & Simplified Conv parameters \\
\texttt{USE\_HW\_WEIGHT\_PRELOAD} & Weight preloading optimization \\
\texttt{USE\_HW\_POOL} & Pooling layer acceleration \\
\bottomrule
\end{tabular}
\end{table}

\begin{table}[htbp]
\caption{Performance Results at Each Optimization Stage}
\label{tab:results}
\centering
\begin{tabular}{@{}lccc@{}}
\toprule
Optimization Stage & Time (ms) & Accuracy & Speedup \\
\midrule
Baseline (Software only) & 54,523 & 98.00\% & 1.00$\times$ \\
+FC Accel (\texttt{USE\_HW\_DOT}) & 52,185 & 98.00\% & 1.04$\times$ \\
+Conv3D (\texttt{USE\_HW\_CONV3D}) & 8,290 & 98.00\% & 6.58$\times$ \\
+Direct Gather & 7,194 & 98.00\% & 7.58$\times$ \\
+Simplified Params & 7,194 & 98.00\% & 7.58$\times$ \\
+Weight Preload & 5,098 & 98.00\% & 10.69$\times$ \\
+Pool Accel (\texttt{USE\_HW\_POOL}) & 4,285 & 98.00\% & 12.72$\times$ \\
\bottomrule
\end{tabular}
\end{table}

\subsection{Stage 1: FC Layer Acceleration }

The first optimization targets the fully connected layer's forward propagation. The FC layer performs matrix-vector multiplication, which can be decomposed into multiple dot product operations.

\textbf{Implementation:} The software calls the DSA to compute dot products between weight rows and input activations. Each output neuron requires one dot product operation.

\textbf{Result:} Execution time reduced from 54,523 ms to 52,185 ms (1.04$\times$ speedup). 

\subsection{Stage 2: Convolution Layer Acceleration }

The convolution operation dominates the computational workload in CNNs. We accelerate the \texttt{conv\_3d()} function by offloading the 3D dot product computation to the DSA.

\textbf{Implementation:} For each output pixel, the software extracts a 3D patch from the input feature map and sends it along with the corresponding kernel weights to the DSA. The DSA computes the dot product and returns the result.

\textbf{Result:} Execution time reduced from 52,185 ms to 8,290 ms (6.58$\times$ speedup compared to baseline). This dramatic improvement confirms that convolution is the primary computational bottleneck.

\subsection{Stage 3: Direct Gather Optimization}

The original Conv3D implementation copies 3D patches to a temporary buffer before sending to DSA, introducing unnecessary memory copy overhead.


\textbf{Result:} Execution time reduced from 8,290 ms to 7,194 ms (7.58$\times$ speedup compared to baseline). This 13\% improvement comes purely from eliminating memory copy overhead.

\subsection{Stage 4: Simplified Convolution Parameters }

Analysis of the CNN model reveals that all convolution layers use padding=0 and stride=1. We exploit this by removing parameter checking and simplifying address calculations.

\textbf{Implementation:} Hardcoded padding and stride values eliminate conditional branches and reduce address computation overhead. The function \texttt{hw\_dot\_conv3d\_patch\_fast()} uses simplified indexing: \texttt{base\_pos = (ox, oy)} directly.

\textbf{Result:} Execution time remains at 7,194 ms. The simplification primarily reduces code complexity and enables further optimizations.

\subsection{Stage 5: Weight Preloading }

Analysis of MMIO traffic reveals that weight writes dominate the communication overhead. For each convolution pixel:
\begin{itemize}
    \item Previous: 75 input values + 75 weight values = 150 MMIO writes
    \item Optimized: 75 input values only = 75 MMIO writes (50\% reduction)
\end{itemize}

\textbf{Implementation:} Weights are preloaded once per output channel using \texttt{hw\_preload\_weights()}. The DSA retains weights in its internal buffer B until explicitly updated. Subsequent dot product operations only need to write input data to buffer A.

\textbf{Result:} Execution time reduced from 7,194 ms to 5,098 ms (10.69$\times$ speedup compared to baseline). This significant improvement confirms that MMIO overhead is the dominant bottleneck.

\subsection{Stage 6: Pooling Layer Acceleration}

Average pooling computes the mean of a 2$\times$2 region, which can be expressed as a dot product with scale factors.

\textbf{Implementation:} The function \texttt{hw\_avg\_pool\_2x2()} sends 4 input values and 4 scale factors (0.25 each) to the DSA. The hardware computes: $result = \sum_{i=0}^{3} input[i] \times 0.25$.

\textbf{Result:} Execution time reduced from 5,098 ms to 4,285 ms (12.72$\times$ speedup compared to baseline). This is our final optimized implementation.

\section{Performance Analysis}

\subsection{Time Breakdown: Data Feeding vs. Computation}

To understand the performance bottleneck, we analyze the time spent on MMIO data transfer (data feeding) versus actual FMA computation. Table~\ref{tab:time_breakdown} shows the estimated time breakdown for different optimization stages.

\begin{table}[htbp]
\caption{Time Breakdown: Data Feeding vs. Computation}
\label{tab:time_breakdown}
\centering
\begin{tabular}{@{}lcccc@{}}
\toprule
\multirow{2}{*}{Stage} & \multicolumn{2}{c}{Time (ms)} & \multicolumn{2}{c}{Percentage} \\
\cmidrule(lr){2-3} \cmidrule(lr){4-5}
 & Feeding & Compute & Feeding & Compute \\
\midrule
Conv3D Accel. & 5,803 & 2,487 & 70\% & 30\% \\
+Direct Gather & 5,036 & 2,158 & 70\% & 30\% \\
+Simplified & 5,036 & 2,158 & 70\% & 30\% \\
+Weight Preload & 2,753 & 2,345 & 54\% & 46\% \\
+Pool Accel. & 2,314 & 1,971 & 54\% & 46\% \\
\bottomrule
\end{tabular}
\end{table}

\subsubsection{Analysis Methodology}
The time breakdown is estimated based on the following parameters:
\begin{itemize}
    \item \textbf{MMIO write latency}: 3-5 cycles per 32-bit write at 41 MHz
    \item \textbf{FMA latency}: 8 cycles pipeline latency per operation
    \item \textbf{Conv2 workload}: $8 \times 8 \times 12$ output pixels $\times$ 75 elements per dot product
\end{itemize}

For a single Conv3D operation with 5$\times$5$\times$3 kernel (75 elements):
\begin{itemize}
    \item \textbf{Data feeding (no preload)}: 75 inputs + 75 weights = 150 MMIO writes $\approx$ 600 cycles
    \item \textbf{Data feeding (with preload)}: 75 inputs only = 75 MMIO writes $\approx$ 300 cycles
    \item \textbf{FMA computation}: 75 FMA operations with 4-way unrolling $\approx$ 19 iterations + pipeline $\approx$ 180 cycles
\end{itemize}

\subsubsection{Key Observation}
The analysis reveals that \textbf{data feeding dominates execution time}, accounting for approximately 70\% of total time before weight preloading optimization. This explains why 4-way loop unrolling (which only accelerates computation) provides diminishing returns.

\subsection{Speedup Breakdown}

Figure~\ref{fig:speedup} illustrates the cumulative speedup achieved at each optimization stage.

\begin{table}[htbp]
\centering
\caption{Cumulative speedup at each optimization stage}
\label{fig:speedup}
\begin{tabular}{|l|r|r|r|}
\hline
\textbf{Stage} & \textbf{Time (ms)} & \textbf{$\Delta$ (ms)} & \textbf{Speedup} \\
\hline
Baseline & 54,523 & -- & 1.00$\times$ \\
+FC Accel & 52,185 & -2,338 & 1.04$\times$ \\
+Conv3D & 8,290 & -43,895 & 6.58$\times$ \\
+Direct Gather & 7,194 & -1,096 & 7.58$\times$ \\
+Simplified & 7,194 & 0 & 7.58$\times$ \\
+Weight Preload & 5,098 & -2,096 & 10.69$\times$ \\
+Pool Accel & 4,285 & -813 & 12.72$\times$ \\
\hline
\end{tabular}
\end{table}

The table shows that Conv3D acceleration contributes the largest improvement (43.9 seconds reduction), confirming that convolution is the dominant computational bottleneck. Weight preloading provides the second-largest improvement by eliminating 50\% of MMIO writes.

\subsection{Bottleneck Analysis}

The performance analysis reveals several insights:

\begin{enumerate}
    \item \textbf{Convolution dominates computation}: The 6.58$\times$ speedup from Conv3D acceleration alone indicates that convolution accounts for approximately 85\% of total execution time in the baseline.
    
    \item \textbf{MMIO overhead is the primary bottleneck}: After hardware acceleration, data feeding via MMIO accounts for approximately 70\% of remaining execution time. The breakdown is:
    \begin{itemize}
        \item Per MMIO write: 4 cycles average at 41 MHz $\approx$ 97 ns
        \item Per Conv3D operation (75 elements, no preload): 150 writes $\approx$ 14.6 $\mu$s for feeding
        \item FMA computation (75 elements): $\approx$ 6.0 $\mu$s
        \item \textbf{Ratio}: Feeding : Compute $\approx$ 70\% : 30\%
    \end{itemize}
    
    \item \textbf{FC layer contribution is minimal}: The FC layers account for less than 5\% of total execution time, making FC acceleration less impactful.
    
    \item \textbf{Weight preloading is highly effective}: By eliminating weight transfers (50\% of MMIO writes), execution time is reduced by 29\% (from 7,194 ms to 5,098 ms), changing the feeding ratio from 70\% to approximately 54\%.
    
    \item \textbf{Pooling acceleration provides additional gains}: Hardware pooling reduces execution time by 16\% (from 5,098 ms to 4,285 ms).
\end{enumerate}

\subsection{Resource Utilization}

The DSA implementation on Arty A7-100T consumes the following FPGA resources:

\begin{itemize}
    \item LUTs: Approximately 35,000 / 63,400 (55\%)
    \item BRAM: 4 / 135 blocks (3\%)
    \item DSP: 3 / 240 slices (1\%)
\end{itemize}

The relatively low DSP utilization indicates room for parallel FMA implementations, though LUT constraints limit the number of concurrent units.


\section{Conclusion}

This work demonstrates effective hardware-software co-design for CNN acceleration on the RISC-V Aquila SoC. Through systematic optimization with six incremental stages, we achieved a 12.72$\times$ speedup while maintaining 98\% classification accuracy. The key findings include:

\begin{enumerate}
    \item \textbf{Convolution acceleration is essential}: Conv3D acceleration provides the most significant improvement (6.58$\times$), confirming that convolution dominates CNN computation.
    
    \item \textbf{MMIO overhead is the limiting factor}: After hardware acceleration, data feeding accounts for 70\% of execution time. Weight preloading reduces this to 54\%.
    
    \item \textbf{Software optimizations are effective}: Direct gather and simplified parameters provide meaningful gains (13\% improvement) with minimal hardware changes.
    
    \item \textbf{Incremental optimization is valuable}: The compile-time flag approach allows systematic evaluation and easy configuration of optimization levels.
\end{enumerate}




\end{document}
